% 图 37.1 —— 由 `checks/emit_hand.py` 自动拼出（probe 精测 + prims2 图元分解）
%%
% ⚠ 这是**稿子不是成品**：跑 `checks/checkhand.py 37.1` 看指标 + 看叠合图。
%%
%% ⚠ 待核：
%   · 本图 probe 报出阴影族 1 个 —— **本脚本不生成阴影**（要照原书手写）；prims2 的 strokes 里可能含阴影线，核对时留意别画重
%   · 可疑字形 1 项（空心圈 ⊙ 常被认成字母 o）—— 见文件末
%%
\documentclass[12pt]{standalone}
\usepackage{fontspec}
\setsansfont{Arial}
\newfontfamily\mfallback{DejaVu Sans}
\usepackage{tikz}
\begin{document}
\begin{tikzpicture}[x=1pt, y=-1pt, line width=4.5pt, line cap=round, line join=round]

  %% ════ 直线段（prims2 的 strokes）════
  \draw[line width=4.5pt] (264.0,192.0) -- (234.0,191.0);
  \draw[line width=5.7pt] (61.0,367.0) -- (58.0,363.0);
  \draw[line width=4.9pt] (318.0,366.0) -- (322.4,364.6);
  \draw[line width=6.0pt] (317.0,367.0) -- (315.0,376.0);
  \draw[line width=4.0pt] (418.0,261.0) -- (406.0,263.0);
  \draw[line width=6.7pt] (313.0,198.0) -- (311.0,192.0);
  \draw[line width=6.3pt] (338.0,523.0) -- (344.0,524.0);
  \draw[line width=5.0pt] (315.0,423.0) -- (315.0,536.0);
  \draw[line width=4.5pt] (173.0,367.0) -- (146.0,368.0);
  \draw[line width=6.0pt] (266.0,192.0) -- (274.0,190.0);
  \draw[line width=4.0pt] (57.0,285.0) -- (57.0,362.0);
  \draw[line width=4.5pt] (229.0,445.0) -- (229.0,422.0);
  \draw[line width=5.0pt] (204.0,366.0) -- (202.0,344.7);
  \draw[line width=5.0pt] (263.6,198.4) -- (265.0,193.0);
  \draw[line width=6.6pt] (345.0,185.0) -- (350.0,189.0);
  \draw[line width=5.0pt] (308.4,423.0) -- (314.0,422.0);
  \draw[line width=5.0pt] (229.0,447.0) -- (229.0,535.0);
  \draw[line width=8.0pt] (58.0,275.0) -- (58.0,284.0);
  \draw[line width=4.2pt] (314.0,537.0) -- (311.0,540.0);
  \draw[line width=3.0pt] (201.0,387.0) -- (188.0,398.0);
  \draw[line width=5.0pt] (142.0,368.0) -- (62.0,368.0);
  \draw[line width=9.0pt] (202.0,395.0) -- (203.0,387.0);
  \draw[line width=7.0pt] (315.0,276.0) -- (315.0,284.0);
  \draw[line width=6.0pt] (224.0,368.0) -- (232.0,362.0);
  \draw[line width=5.0pt] (224.0,368.0) -- (208.0,368.0);
  \draw[line width=7.0pt] (224.0,368.0) -- (228.0,373.0);
  \draw[line width=6.6pt] (349.0,280.0) -- (344.0,276.0);
  \draw[line width=7.0pt] (62.0,192.0) -- (58.0,187.0);
  \draw[line width=5.0pt] (62.0,192.0) -- (182.0,192.0);
  \draw[line width=8.0pt] (62.0,192.0) -- (57.0,197.0);
  \draw[line width=6.7pt] (316.0,537.0) -- (321.0,540.0);
  \draw[line width=4.0pt] (316.0,275.0) -- (314.0,201.0);
  \draw[line width=6.3pt] (338.0,521.0) -- (337.0,515.0);
  \draw[line width=4.3pt] (315.0,421.0) -- (316.0,417.0);
  \draw[line width=9.4pt] (203.0,387.0) -- (211.0,383.0);
  \draw[line width=5.9pt] (334.0,275.0) -- (333.0,269.4);
  \draw[line width=5.9pt] (340.0,190.6) -- (341.0,185.0);
  \draw[line width=8.0pt] (334.0,275.0) -- (342.0,275.0);
  \draw[line width=5.1pt] (334.0,275.0) -- (330.0,278.0);
  \draw[line width=4.0pt] (176.0,369.0) -- (186.0,397.0);
  \draw[line width=5.0pt] (453.0,169.0) -- (452.0,174.0);
  \draw[line width=6.0pt] (228.0,536.0) -- (224.0,540.0);
  \draw[line width=4.0pt] (394.0,172.0) -- (365.0,179.0);
  \draw[line width=3.8pt] (317.0,276.0) -- (320.0,278.0);
  \draw[line width=4.0pt] (310.0,191.0) -- (274.0,190.0);
  \draw[line width=4.0pt] (363.0,179.0) -- (346.0,184.0);
  \draw[line width=6.7pt] (345.0,183.0) -- (347.0,177.0);
  \draw[line width=5.0pt] (317.0,302.0) -- (316.0,285.0);
  \draw[line width=5.0pt] (229.0,392.0) -- (229.0,374.0);
  \draw[line width=5.3pt] (183.0,192.0) -- (188.0,191.0);
  \draw[line width=5.0pt] (228.0,393.0) -- (211.0,383.0);
  \draw[line width=5.0pt] (315.0,321.0) -- (317.0,366.0);
  \draw[line width=5.5pt] (337.0,523.0) -- (325.0,532.0);
  \draw[line width=6.1pt] (234.0,540.0) -- (230.0,536.0);
  \draw[line width=4.0pt] (229.0,420.0) -- (229.0,395.0);
  \draw[line width=3.0pt] (122.0,440.0) -- (164.0,412.0);
  \draw[line width=4.0pt] (433.0,163.0) -- (397.0,171.0);
  \draw[line width=3.0pt] (344.0,185.0) -- (341.0,185.0);
  \draw[line width=5.0pt] (176.0,368.0) -- (204.0,367.0);
  \draw[line width=4.0pt] (211.0,382.0) -- (207.0,369.0);
  \draw[line width=4.0pt] (315.0,377.0) -- (316.0,415.0);
  \draw[line width=5.0pt] (308.0,326.0) -- (314.0,321.0);
  \draw[line width=4.0pt] (371.0,495.0) -- (339.0,521.0);
  \draw[line width=5.0pt] (328.0,187.0) -- (338.0,184.0);
  \draw[line width=4.0pt] (322.0,541.0) -- (558.7,542.0);
  \draw[line width=5.0pt] (558.7,542.0) -- (567.6,539.4);
  \draw[line width=4.0pt] (567.6,539.4) -- (567.7,16.7);
  \draw[line width=4.0pt] (567.7,16.7) -- (58.6,16.4);
  \draw[line width=4.0pt] (58.6,16.4) -- (57.0,186.0);
  \draw[line width=5.0pt] (57.0,197.0) -- (57.0,274.0);
  \draw[line width=5.0pt] (226.0,192.0) -- (190.0,191.0);
  \draw[line width=6.7pt] (57.0,373.0) -- (61.0,369.0);
  \draw[line width=4.0pt] (57.0,373.0) -- (57.4,540.4);
  \draw[line width=4.0pt] (57.4,540.4) -- (223.0,541.0);
  \draw[line width=4.0pt] (345.0,275.0) -- (365.0,271.0);
  \draw[line width=5.3pt] (228.0,192.0) -- (233.0,191.0);
  \draw[line width=4.0pt] (167.0,412.0) -- (186.0,398.0);
  \draw[line width=3.4pt] (341.0,185.0) -- (338.0,184.0);
  \draw[line width=4.0pt] (235.0,541.0) -- (310.0,541.0);
  \draw[line width=5.0pt] (343.0,274.0) -- (347.0,267.0);
  \draw[line width=5.0pt] (313.0,201.0) -- (307.3,207.7);
  \draw[line width=5.0pt] (307.3,207.7) -- (233.0,361.0);
  \draw[line width=6.0pt] (315.0,320.0) -- (317.0,303.0);
  \draw[line width=4.0pt] (214.1,158.9) -- (234.1,141.9);
  \draw[line width=5.0pt] (329.0,278.0) -- (322.0,278.0);
  \draw[line width=4.0pt] (369.0,269.0) -- (403.0,264.0);
  \draw[line width=7.1pt] (321.0,279.0) -- (317.0,284.0);

  %% ════ 曲线（prims2 的 curves —— 三段式，坐标同为 (y,x)）════
  \draw[line width=5.0pt] (405.0,262.0) .. controls (408.3,231.4) and (408.5,199.3) .. (396.0,172.0);
  \draw[line width=4.0pt] (322.4,364.6) .. controls (344.3,337.3) and (359.4,305.6) .. (366.0,272.0);
  \draw[line width=6.0pt] (224.0,542.0) .. controls (224.8,547.1) and (234.5,548.0) .. (234.0,542.0);
  \draw[line width=4.0pt] (318.0,302.0) .. controls (326.0,297.0) and (326.9,286.9) .. (330.0,279.0);
  \draw[line width=5.0pt] (182.0,193.0) .. controls (150.1,244.1) and (134.3,305.5) .. (143.0,367.0);
  \draw[line width=5.0pt] (368.0,268.0) .. controls (370.2,238.3) and (375.4,208.1) .. (365.0,180.0);
  \draw[line width=4.0pt] (174.0,366.0) .. controls (168.4,302.2) and (189.8,242.5) .. (227.0,193.0);
  \draw[line width=4.0pt] (228.0,446.0) .. controls (205.1,442.5) and (180.5,433.1) .. (167.0,412.0);
  \draw[line width=7.0pt] (56.0,275.0) .. controls (50.7,275.8) and (51.3,282.7) .. (56.0,284.0);
  \draw[line width=8.0pt] (315.0,198.0) .. controls (322.5,194.4) and (317.8,182.8) .. (311.0,190.0);
  \draw[line width=5.0pt] (202.0,344.7) .. controls (210.9,292.2) and (228.9,240.8) .. (263.6,198.4);
  \draw[line width=5.0pt] (145.0,369.0) .. controls (146.3,384.9) and (156.2,397.6) .. (164.0,411.0);
  \draw[line width=5.0pt] (230.0,446.0) .. controls (259.0,447.6) and (285.4,438.8) .. (308.4,423.0);
  \draw[line width=4.0pt] (234.0,190.0) .. controls (257.5,144.4) and (340.9,120.4) .. (363.0,179.0);
  \draw[line width=4.0pt] (228.0,421.0) .. controls (212.3,419.9) and (196.1,412.9) .. (188.0,398.0);
  \draw[line width=5.0pt] (29.0,208.0) .. controls (36.6,217.2) and (38.0,193.9) .. (30.0,199.0);
  \draw[line width=5.0pt] (452.0,168.0) .. controls (443.6,171.0) and (442.5,158.2) .. (449.0,156.0);
  \draw[line width=6.0pt] (449.0,156.0) .. controls (458.0,153.7) and (455.8,163.4) .. (453.0,168.0);
  \draw[line width=4.0pt] (333.0,269.4) .. controls (340.9,245.0) and (347.1,216.8) .. (340.0,190.6);
  \draw[line width=7.0pt] (56.0,363.0) .. controls (51.5,365.4) and (52.2,371.7) .. (57.0,373.0);
  \draw[line width=5.0pt] (230.0,421.0) .. controls (264.1,422.4) and (288.6,396.5) .. (314.0,377.0);
  \draw[line width=5.0pt] (230.0,394.0) .. controls (265.7,389.6) and (291.1,355.1) .. (308.0,326.0);
  \draw[line width=7.0pt] (321.0,542.0) .. controls (320.4,547.3) and (311.6,547.3) .. (311.0,542.0);
  \draw[line width=7.0pt] (57.0,197.0) .. controls (51.3,196.3) and (50.8,189.0) .. (56.0,187.0);
  \draw[line width=5.0pt] (274.0,190.0) .. controls (285.6,169.9) and (326.1,155.1) .. (338.0,184.0);
  \draw[line width=6.0pt] (230.0,373.0) .. controls (235.6,373.3) and (236.1,365.9) .. (233.0,363.0);
  \draw[line width=4.0pt] (189.0,190.0) .. controls (192.7,177.4) and (205.4,168.4) .. (214.1,158.9);
  \draw[line width=4.0pt] (234.1,141.9) .. controls (281.4,98.4) and (369.2,107.9) .. (395.0,171.0);
  \draw[line width=5.0pt] (317.0,416.0) .. controls (366.9,380.7) and (395.7,322.8) .. (404.0,265.0);

  %% ════ 标签（probe 直接给的 \node 行）════
  %% ⚠ 全部补了 fill=white：文字在最上层，否则与曲线交叉干扰
     \node[anchor=center,inner sep=0pt,fill=white,font=\fontsize{24.7}{30.9}\selectfont] at (33.6,179.9) {\textsf{\textbf{2}}};   % box(27,171,12,17)
     \node[anchor=center,inner sep=0pt,fill=white,font=\fontsize{30.8}{38.5}\selectfont] at (41.4,197.1) {{\mfallback\textbf{−}}};   % box(27,190,14,5)
     \node[anchor=center,inner sep=0pt,fill=white,font=\fontsize{33.8}{42.3}\selectfont] at (479.1,261.5) {\textsf{\textbf{×}}};   % box(471,251,15,16)
     \node[anchor=center,inner sep=0pt,fill=white,font=\fontsize{28.3}{35.4}\selectfont] at (510.4,262.4) {\textsf{\textbf{\textit{x}}}};   % box(502,252,15,17)
     \node[anchor=center,inner sep=0pt,fill=white,font=\fontsize{27.4}{34.3}\selectfont] at (439.4,262.8) {\textsf{\textbf{\textit{x}}}};   % box(431,253,15,16)
     \node[anchor=center,inner sep=0pt,fill=white,font=\fontsize{23.7}{29.7}\selectfont] at (454.8,271.4) {\textsf{\textbf{1}}};   % box(448,263,8,16)
     \node[anchor=center,inner sep=0pt,fill=white,font=\fontsize{23.6}{29.6}\selectfont] at (524.1,271.9) {\textsf{\textbf{2}}};   % box(518,263,11,17)
     \node[anchor=center,inner sep=0pt,fill=white,font=\fontsize{27.4}{34.3}\selectfont] at (21.4,279.8) {\textsf{\textbf{\textit{x}}}};   % box(13,270,15,16)
     \node[anchor=center,inner sep=0pt,fill=white,font=\fontsize{23.6}{29.6}\selectfont] at (35.1,288.9) {\textsf{\textbf{2}}};   % box(29,280,11,17)
     \node[anchor=center,inner sep=0pt,fill=white,font=\fontsize{21.5}{26.9}\selectfont] at (36.0,359.9) {\textsf{\textbf{1}}};   % box(30,352,7,15)
     \node[anchor=center,inner sep=0pt,fill=white,font=\fontsize{31.9}{39.8}\selectfont] at (41.2,375.3) {{\mfallback\textbf{−}}};   % box(26,368,15,5)
     \node[anchor=center,inner sep=0pt,fill=white,font=\fontsize{22.3}{27.9}\selectfont] at (33.9,385.6) {\textsf{\textbf{3}}};   % box(28,377,11,16)
     \node[anchor=center,inner sep=0pt,fill=white,font=\fontsize{30.0}{37.4}\selectfont] at (108.5,453.0) {\textsf{\textbf{\textit{P}}}};   % box(97,441,19,22)
     \node[anchor=center,inner sep=0pt,fill=white,font=\fontsize{27.4}{34.3}\selectfont] at (388.4,492.8) {\textsf{\textbf{\textit{x}}}};   % box(380,483,15,16)
     \node[anchor=center,inner sep=0pt,fill=white,font=\fontsize{20.6}{25.7}\selectfont] at (403.4,503.3) {\textsf{\textbf{1}}};   % box(398,495,6,16)
     \node[anchor=center,inner sep=0pt,fill=white,font=\fontsize{22.2}{27.8}\selectfont] at (232.1,570.4) {\textsf{\textbf{1}}};   % box(226,562,7,16)
     \node[anchor=center,inner sep=0pt,fill=white,font=\fontsize{21.5}{26.9}\selectfont] at (318.0,570.9) {\textsf{\textbf{1}}};   % box(312,563,7,15)
     \node[anchor=center,inner sep=0pt,fill=white,font=\fontsize{30.8}{38.5}\selectfont] at (236.4,586.1) {{\mfallback\textbf{−}}};   % box(222,579,14,5)
     \node[anchor=center,inner sep=0pt,fill=white,font=\fontsize{27.5}{34.4}\selectfont] at (322.7,586.1) {{\mfallback\textbf{−}}};   % box(309,580,14,4)
     \node[anchor=center,inner sep=0pt,fill=white,font=\fontsize{23.0}{28.7}\selectfont] at (229.9,596.1) {\textsf{\textbf{3}}};   % box(224,587,11,17)
     \node[anchor=center,inner sep=0pt,fill=white,font=\fontsize{23.6}{29.6}\selectfont] at (317.1,595.9) {\textsf{\textbf{2}}};   % box(311,587,11,17)

  % 钉死画布：否则超出的图元/控制点会把页面撑大，1:1 回比就失准
  \pgfresetboundingbox
  \path[use as bounding box] (0,0) rectangle (578,612);
\end{tikzpicture}
\end{document}

% ⚠ 待核清单（probe 拿不准，人眼过一遍）：
%   · 未识别/跳过：⑤ 跳过/未识别的字形
